% Preamble / 宏包区
\documentclass{article}
\usepackage[margin=0mm]{geometry}
\usepackage{xcolor}
\usepackage{graphicx}
\usepackage{tikz}
\usepackage{amsmath}

\begin{document}
% Metadata / 注释区
% Exported by guiLaTeX ExportCore
% Export date: 2026-10-01 16:48:02
% Page size: A4 (210mm × 297mm)
% IR Metadata for roundtrip:
% BEGIN_IR_METADATA
% {
%   "export_date": "2026-10-01 16:48:02",
%   "page_size": "A4 (210mm × 297mm)",
%   "exporter": "guiLaTeX ExportCore",
%   "version": "1.0",
%   "ir": {
%     "elements": [
%       {
%         "id": "1",
%         "type": "title",
%         "content": "标题：可视化编辑 LaTeX 文档 (文本框)",
%         "page": 1,
%         "x": 100,
%         "y": 100,
%         "width": 400,
%         "height": 50,
%         "rotation": 0,
%         "layer": 1,
%         "font_family_zh": "Noto Sans SC",
%         "font_family_en": "Inter",
%         "font_size": 24,
%         "color": "#000000",
%         "alignment": "left",
%         "visible": true
%       },
%       {
%         "id": "2",
%         "type": "author",
%         "content": "作者：某某 / 博士生原型测试 (文本框)",
%         "page": 1,
%         "x": 100,
%         "y": 150,
%         "width": 200,
%         "height": 30,
%         "rotation": 0,
%         "layer": 2,
%         "font_family_zh": "Noto Sans SC",
%         "font_family_en": "Inter",
%         "font_size": 16,
%         "color": "#000000",
%         "alignment": "left",
%         "visible": true
%       },
%       {
%         "id": "3",
%         "type": "paragraph",
%         "content": "这是一段用于演示段落对象的正文内容，可用于测试拖动、缩放与排版差异。(段落)",
%         "page": 1,
%         "x": 100,
%         "y": 200,
%         "width": 400,
%         "height": 80,
%         "rotation": 0,
%         "layer": 3,
%         "font_family_zh": "Noto Sans SC",
%         "font_family_en": "Inter",
%         "font_size": 14,
%         "color": "#000000",
%         "alignment": "left",
%         "visible": true
%       },
%       {
%         "id": "4",
%         "type": "equation",
%         "content": "E = mc^2",
%         "page": 1,
%         "x": 100,
%         "y": 280,
%         "width": 150,
%         "height": 40,
%         "rotation": 0,
%         "layer": 4,
%         "font_family_zh": "Noto Sans SC",
%         "font_family_en": "Inter",
%         "font_size": 18,
%         "color": "#000000",
%         "alignment": "left",
%         "visible": true
%       },
%       {
%         "id": "5",
%         "type": "image",
%         "content": "图片示例 / 占位图",
%         "page": 1,
%         "x": 100,
%         "y": 330,
%         "width": 200,
%         "height": 150,
%         "rotation": 0,
%         "layer": 5,
%         "font_family_zh": "SimSun",
%         "font_family_en": "Times New Roman",
%         "font_size": 12,
%         "color": "#000000",
%         "alignment": "left",
%         "visible": true
%       }
%     ]
%   }
% }
% END_IR_METADATA
% Optional semantic summary 区
% Title: 标题：可视化编辑 LaTeX 文档 (文本框)
% Author: 作者：某某 / 博士生原型测试 (文本框)
% Absolute positioned objects 区
\begin{tikzpicture}[remember picture, overlay]
  % Layer 1: 1
  \node[anchor=north west, rotate=0, text width=400mm, align=left, font=\fontsize{24}{28.8}\selectfont, color={rgb,1:red,0.0;green,0.0;blue,0.0}, visible on layer=1]
    at (100mm, 197mm) {标题：可视化编辑 LaTeX 文档 (文本框)};
  % Layer 2: 2
  \node[anchor=north west, rotate=0, text width=200mm, align=left, font=\fontsize{16}{19.2}\selectfont, color={rgb,1:red,0.0;green,0.0;blue,0.0}, visible on layer=2]
    at (100mm, 147mm) {作者：某某 / 博士生原型测试 (文本框)};
  % Layer 3: 3
  \node[anchor=north west, rotate=0, text width=400mm, align=left, font=\fontsize{14}{16.8}\selectfont, color={rgb,1:red,0.0;green,0.0;blue,0.0}, visible on layer=3]
    at (100mm, 97mm) {这是一段用于演示段落对象的正文内容，可用于测试拖动、缩放与排版差异。(段落)};
  % Layer 4: 4
  \node[anchor=north west, rotate=0, text width=150mm, align=left, font=\fontsize{18}{21.6}\selectfont, color={rgb,1:red,0.0;green,0.0;blue,0.0}, visible on layer=4]
    at (100mm, 17mm) {$E = mc^2$};
  % Layer 5: 5
  \node[anchor=north west, rotate=0, text width=200mm, align=left, font=\fontsize{12}{14.4}\selectfont, color={rgb,1:red,0.0;green,0.0;blue,0.0}, visible on layer=5]
    at (100mm, -33mm) {\includegraphics[width=200mm, height=150mm]{图片示例 / 占位图}};
\end{tikzpicture}
\end{document}